\documentclass[12pt,a4paper]{article}

\usepackage[margin=2.2cm]{geometry}
\usepackage{setspace}
\usepackage{titlesec}
\usepackage{tocloft}
\usepackage{graphicx}
\usepackage{float}
\usepackage{booktabs}
\usepackage{longtable}
\usepackage{array}
\usepackage{amsmath}
\usepackage{amssymb}
\usepackage{hyperref}
\usepackage{xcolor}
\usepackage{listings}
\usepackage{caption}
\usepackage{iftex}

\ifPDFTeX
  \usepackage[T5]{fontenc}
  \usepackage[utf8]{inputenc}
\else
  \usepackage{fontspec}
  \setmainfont{Times New Roman}
\fi

\onehalfspacing

\hypersetup{
  colorlinks=true,
  linkcolor=blue!50!black,
  urlcolor=blue!50!black,
  citecolor=blue!50!black
}

\titleformat{\section}{\large\bfseries}{\thesection.}{0.5em}{}
\titleformat{\subsection}{\normalsize\bfseries}{\thesubsection.}{0.5em}{}
\titleformat{\subsubsection}{\normalsize\bfseries}{\thesubsubsection.}{0.5em}{}

\lstdefinestyle{pythonstyle}{
  language=Python,
  basicstyle=\ttfamily\small,
  keywordstyle=\color{blue!70!black}\bfseries,
  commentstyle=\color{green!40!black},
  stringstyle=\color{red!50!black},
  numbers=left,
  numberstyle=\tiny\color{gray},
  stepnumber=1,
  numbersep=8pt,
  showstringspaces=false,
  breaklines=true,
  frame=single,
  tabsize=2
}

\lstdefinestyle{yamlstyle}{
  basicstyle=\ttfamily\small,
  numbers=left,
  numberstyle=\tiny\color{gray},
  stepnumber=1,
  numbersep=8pt,
  showstringspaces=false,
  breaklines=true,
  frame=single,
  tabsize=2
}

\newcommand{\screenshotplaceholder}[2]{
  \begin{figure}[H]
    \centering
    \fbox{\parbox[c][6cm][c]{0.9\textwidth}{\centering \textbf{[Placeholder Screenshot]}\\#1\\\vspace{0.3cm}Chèn ảnh chụp màn hình tại đây.}}
    \caption{#1}
    \label{#2}
  \end{figure}
}

\begin{document}

\tableofcontents
\newpage

\section{Tóm tắt báo cáo}
Báo cáo này trình bày đầy đủ quá trình thực hiện bài tập gồm 6 nhóm yêu cầu: (1) khảo sát các hướng áp dụng AI trong e-commerce, (2) xây dựng ứng dụng phân tích hành vi khách hàng để tư vấn, (3) xây dựng model\_behavior dựa trên Deep Learning, (4) xây dựng Knowledge Base (KB) cho tư vấn, (5) áp dụng Retrieval-Augmented Generation (RAG) cho chatbot, và (6) deploy và tích hợp vào hệ thống e-commerce đã có.

Hệ thống được triển khai theo hướng microservices với Django, kết nối giữa các service thông qua Docker Compose network. Chatbot được tách thành service riêng để đảm bảo đóng gói logic AI, trong khi customer\_service đóng vai trò gateway cho frontend. Quá trình huấn luyện model và tạo dữ liệu train thực tế được tự động hóa bằng command và artifact JSON để dễ kiểm tra, đối chiếu, báo cáo.

\section{Khảo sát các áp dụng AI trong e-commerce}

\subsection{Tổng quan xu hướng}
AI trong e-commerce đã chuyển từ vai trò hỗ trợ sang vai trò cốt lõi trong vận hành và tăng trưởng doanh thu. Các hướng ứng dụng chính:
\begin{itemize}
  \item \textbf{Hệ thống gợi ý sản phẩm (Recommendation Systems):} cá nhân hóa theo lịch sử xem, giỏ hàng, đơn hàng và hành vi click.
  \item \textbf{Chatbot và trợ lý tư vấn thông minh:} tư vấn sản phẩm 24/7, giảm tải cho CSKH, tăng conversion.
  \item \textbf{Dự báo nhu cầu và tồn kho:} sử dụng time-series, causal signal, và event stream.
  \item \textbf{Phát hiện gian lận giao dịch:} anomaly detection và graph-based risk scoring.
  \item \textbf{Tự động hóa nội dung (AIGC):} tạo mô tả sản phẩm, tiêu đề, FAQ, phản hồi email.
  \item \textbf{Search thông minh:} semantic search, vector search, hybrid ranking.
\end{itemize}

\subsection{Bài toán thực tế và tác động kinh doanh}
Trong e-commerce, 3 KPI được theo dõi sát là:
\begin{itemize}
  \item \textbf{Conversion Rate (CR):} tỷ lệ khách hàng thực hiện mua.
  \item \textbf{Average Order Value (AOV):} giá trị đơn hàng trung bình.
  \item \textbf{Retention/Repeat Rate:} tỷ lệ mua lại.
\end{itemize}

AI có tác động trực tiếp qua:
\begin{itemize}
  \item Nhanh hơn trong việc tìm sản phẩm phù hợp.
  \item Đề xuất cross-sell/up-sell đúng ngữ cảnh.
  \item Giảm giám sát và thao tác thủ công của đội vận hành.
\end{itemize}

\subsection{Khoảng trống nghiên cứu và hướng đề tài}
Nhiều hệ thống chatbot chỉ dùng rule-based hoặc keyword. Hướng đề tài này giải quyết khoảng trống bằng cách:
\begin{itemize}
  \item Kết hợp \textbf{phân tích hành vi} + \textbf{RAG} + \textbf{gợi ý sản phẩm}.
  \item Triển khai được trên hệ thống microservices thực tế, không dùng mô hình demo tách rời.
  \item Có cơ chế fallback khi LLM unavailable để đảm bảo tính ổn định vận hành.
\end{itemize}

\screenshotplaceholder{Tổng quan kiến trúc AI trong e-commerce (chèn sơ đồ từ draw.io/mermaid)}{fig:survey-architecture}

\section{Mô tả hệ thống và bài toán dự án}

\subsection{Mục tiêu bài toán}
Xây dựng ứng dụng \textit{Phân tích hành vi khách hàng để tư vấn dịch vụ} trong bối cảnh e-commerce đã có nhiều catalog service (laptop, mobile, accessory), có giao diện customer và staff.

\subsection{Kiến trúc microservices}
Hệ thống gồm các service chính:
\begin{itemize}
  \item customer\_service: giao diện khách hàng, giỏ hàng, checkout, pay, gateway chatbot.
  \item staff\_service: dashboard quản trị và CRUD sản phẩm.
  \item laptop\_service, mobile\_service, accessory\_service: catalog API.
  \item chatbot\_service: behavior AI, KB, RAG, gọi LLM, fallback.
\end{itemize}

\subsection{Trích đoạn cấu hình deploy}
\begin{lstlisting}[style=yamlstyle,caption={Trích đoạn docker-compose cho chatbot và customer integration},label={lst:compose}]
chatbot_service:
  build: ./services/chatbot_service
  hostname: chatbot-service
  environment:
    LAPTOP_SERVICE_URL: http://laptop-service:8000
    MOBILE_SERVICE_URL: http://mobile-service:8000
    ACCESSORY_SERVICE_URL: http://accessory-service:8000
  ports:
    - "8005:8000"

customer_service:
  build: ./services/customer_service
  environment:
    CHATBOT_SERVICE_URL: http://chatbot-service:8000
  depends_on:
    - chatbot_service
\end{lstlisting}

\screenshotplaceholder{Docker Compose containers đang running (docker compose ps)}{fig:docker-ps}

\section{Xây dựng ứng dụng phân tích hành vi khách hàng để tư vấn}

\subsection{Dòng dữ liệu hành vi}
Dòng dữ liệu hành vi được ghi nhận theo user\_ref, current\_product, user\_context (cart/saved/recent\_paid), message và thời gian event.

\subsection{Model dữ liệu hành vi}
Bảng sự kiện hành vi được định nghĩa trong chatbot\_service:
\begin{lstlisting}[style=pythonstyle,caption={Model BehaviorEvent},label={lst:behavior-model}]
class BehaviorEvent(models.Model):
    EVENT_CHATBOT_ASK = "chatbot_ask"
    EVENT_CHOICES = [(EVENT_CHATBOT_ASK, "Chatbot ask")]

    user_ref = models.CharField(max_length=120, db_index=True)
    event_type = models.CharField(max_length=40, choices=EVENT_CHOICES)
    product_service = models.CharField(max_length=20, blank=True)
    product_id = models.PositiveIntegerField(default=0)
    metadata = models.JSONField(default=dict, blank=True)
    created_at = models.DateTimeField(auto_now_add=True)
\end{lstlisting}

\subsection{Ghi sự kiện hành vi trong chatbot pipeline}
\begin{lstlisting}[style=pythonstyle,caption={Hàm ghi sự kiện hành vi khi có câu hỏi chatbot},label={lst:record-behavior}]
def record_behavior_event(user_ref, message, current_product=None, user_context=None):
    BehaviorEvent.objects.create(
        user_ref=user_ref,
        event_type=BehaviorEvent.EVENT_CHATBOT_ASK,
        product_service=str(current_product.get("service") or "")[:20],
        product_id=max(0, int(current_product.get("id") or 0)),
        metadata={
            "message": str(message or "")[:500],
            "current_product_service": str(current_product.get("service") or "")[:20],
            "cart_items_count": len(user_context.get("cart_items") or []),
            "saved_items_count": len(user_context.get("saved_items") or []),
            "recent_paid_items_count": len(user_context.get("recent_paid_items") or []),
        },
    )
\end{lstlisting}

\subsection{Sinh dữ liệu train thực tế từ giao dịch đã có}
Khác với dữ liệu synthetic, dữ liệu train trong dự án được backfill từ order/cart/saved thực tế từ customer\_service.

\begin{lstlisting}[style=pythonstyle,caption={Command backfill dữ liệu hành vi thực tế},label={lst:backfill-cmd}]
python manage.py backfill_chatbot_behavior \
  --max-events 1200 \
  --max-users 300 \
  --source-status paid
\end{lstlisting}

Kết quả vận hành ghi nhận:
\begin{itemize}
  \item sent = 1200
  \item success = 1200
  \item failed = 0
\end{itemize}

\screenshotplaceholder{Kết quả chạy command backfill\_chatbot\_behavior thành công}{fig:backfill-result}

\section{Xây dựng model\_behavior dựa trên Deep Learning}

\subsection{Định nghĩa đặc trưng (features)}
Model sử dụng vector đặc trưng gồm 12 thành phần:
\begin{enumerate}
  \item ask\_count
  \item mentions\_laptop
  \item mentions\_mobile
  \item mentions\_accessory
  \item current\_product\_laptop
  \item current\_product\_mobile
  \item current\_product\_accessory
  \item cart\_items\_count
  \item saved\_items\_count
  \item recent\_paid\_items\_count
  \item avg\_phrase\_len
  \item english\_ratio
\end{enumerate}

Dữ liệu thô được biến đổi bằng log-scale:
\[
\phi(x) = \log(1 + x)
\]

\subsection{Kiến trúc Deep Learning}
Model là MLP (Multi-layer Perceptron) tự xây dựng:
\begin{itemize}
  \item Input: 12
  \item Hidden layer 1: 18 (ReLU)
  \item Hidden layer 2: 10 (ReLU)
  \item Output: 4
\end{itemize}

Trong đó output gồm:
\begin{itemize}
  \item 1 nút sigmoid cho \textit{intent score}
  \item 3 nút softmax cho \textit{category} (laptop/mobile/accessory)
\end{itemize}

\subsection{Công thức học}
\textbf{Sigmoid} cho intent:
\[
\hat{y}_{intent} = \sigma(z) = \frac{1}{1 + e^{-z}}
\]

\textbf{Softmax} cho phân loại danh mục:
\[
\hat{y}_k = \frac{e^{z_k}}{\sum_j e^{z_j}}
\]

\textbf{Hàm mất mát tổng hợp}:
\[
\mathcal{L} = \underbrace{-\left[y\log(\hat{y}) + (1-y)\log(1-\hat{y})\right]}_{\text{BCE cho intent}} +
\underbrace{-\log(\hat{y}_{c})}_{\text{CE cho category}}
\]

\subsection{Trích đoạn code huấn luyện}
\begin{lstlisting}[style=pythonstyle,caption={Train model\_behavior và lưu artifact},label={lst:train-behavior}]
def train_and_save_behavior_model(epochs=120, lr=0.02):
    dataset_bundle = _build_training_dataset()
    dataset = dataset_bundle.get("augmented_dataset") or []
    model, metrics = _fit_model(dataset, epochs=epochs, lr=lr)

    payload = {
        "feature_names": FEATURE_NAMES,
        "service_names": SERVICE_NAMES,
        "model": model,
        "metrics": metrics,
        "version": 1,
    }

    MODEL_BEHAVIOR_PATH.write_text(json.dumps(payload), encoding="utf-8")
    _save_training_data_snapshot(dataset_bundle)
    return payload
\end{lstlisting}

\subsection{Kết quả huấn luyện}
Sau khi backfill dữ liệu thực tế và train:
\begin{itemize}
  \item samples = 3600
  \item loss = 0.009914
\end{itemize}

Artifact được lưu:
\begin{itemize}
  \item model weights: \texttt{model\_behavior.json}
  \item training snapshot dễ đọc: \texttt{training\_data\_behavior.json}
\end{itemize}

\screenshotplaceholder{File model\_behavior.json và metrics sau train}{fig:model-artifact}

\section{Xây dựng Knowledge Base (KB) cho tư vấn}

\subsection{Nguồn dữ liệu KB}
KB được tổng hợp từ:
\begin{itemize}
  \item FAQ nội bộ (shipping, delivery, return)
  \item Product catalog từ 3 service (laptop/mobile/accessory)
\end{itemize}

\subsection{Pipeline tạo KB}
\begin{lstlisting}[style=pythonstyle,caption={Build KB artifact},label={lst:build-kb}]
def build_and_save_knowledge_base(max_products_per_service=120):
    docs = []
    docs.extend(_faq_docs())

    for service_name, base_url in _service_urls().items():
        docs.extend(_fetch_products(service_name, base_url, limit=max_products_per_service))

    for doc in docs:
        doc["tokens"] = _tokenize(" ".join([doc.get("title") or "", doc.get("text") or ""]))

    payload = {"version": 1, "documents": docs, "stats": {...}}
    KB_PATH.write_text(json.dumps(payload), encoding="utf-8")
\end{lstlisting}

\subsection{Cơ cấu document trong KB}
Mỗi tài liệu trong KB có các trường:
\begin{itemize}
  \item doc\_id, doc\_type
  \item service, product\_id
  \item title, text, url
  \item tokens
\end{itemize}

\screenshotplaceholder{Nội dung file knowledge\_base.json sau khi build}{fig:kb-artifact}

\section{Áp dụng RAG để xây dựng chat tư vấn}

\subsection{Mục tiêu RAG}
Thay vì cho LLM trả lời kiến thức chung, RAG bổ sung ngữ cảnh retrieval từ KB và candidate sản phẩm thực tế, giúp:
\begin{itemize}
  \item Giảm hallucination
  \item Tăng tính truy vết (có citations)
  \item Tăng khả năng gợi ý có căn cứ vào inventory thực
\end{itemize}

\subsection{Luồng RAG trong chatbot\_service}
\begin{enumerate}
  \item Nhận câu hỏi + user context + current product
  \item Phân tích behavior signal
  \item Retrieve top-k tài liệu từ KB
  \item Xây prompt có: question + profile + behavior + rag docs + recommendation candidates
  \item Gọi LLM (Gemma/Gemini)
  \item Nếu lỗi: fallback rule-based
  \item Trả answer + recommendations + citations
\end{enumerate}

\subsection{Trích đoạn retrieval và citation}
\begin{lstlisting}[style=pythonstyle,caption={Retrieve RAG context và tạo citation},label={lst:rag-retrieve}]
def retrieve_rag_context(question, preferred_services=None, current_product=None, top_k=6):
    payload = load_knowledge_base(auto_build=True)
    docs = payload.get("documents") or []
    ...
    ranked.sort(key=lambda row: row[0], reverse=True)
    return selected[:top_k]


def rag_citations_from_docs(docs, limit=3):
    citations = []
    for doc in docs:
        ...
        citations.append({"label": label, "detail": detail, "url": url})
    return citations
\end{lstlisting}

\subsection{Trích đoạn generation pipeline}
\begin{lstlisting}[style=pythonstyle,caption={Pipeline tạo chatbot response có RAG và fallback},label={lst:chatbot-pipeline}]
def generate_chatbot_response(question, current_product=None, user_context=None, user_ref="", limit=5):
    record_behavior_event(...)
    behavior_signal = predict_behavior_for_user_ref(...)
    recommendations = recommend_products(...)
    rag_docs = retrieve_rag_context(...)

    prompt_text = _build_prompt(..., rag_docs=rag_docs, ...)
    llm_answer, error_code, llm_source = _call_llm(prompt_text, max_output_tokens=320)

    if llm_answer:
        answer = _build_focused_answer(llm_answer, recommendations, language)
    else:
        answer = _fallback_answer(recommendations, language, error_code=error_code)

    citations = rag_citations_from_docs(rag_docs, limit=3)
    return {"answer": answer, "recommendations": recommendations, "citations": citations, ...}
\end{lstlisting}

\subsection{Fallback và tính sẵn sàng vận hành}
Nếu LLM không sẵn sàng (missing key, timeout, 429, network error), hệ thống vẫn trả lời được bằng fallback mode. Đây là điểm quan trọng trong hệ thống production-like vì đảm bảo UX liên tục, không bị treo.

\screenshotplaceholder{Giao diện chat hiển thị answer + recommendation + citations}{fig:chat-ui}

\section{Deploy và tích hợp trong hệ thống e-commerce}

\subsection{Tích hợp customer gateway -> chatbot service}
Frontend gửi tới endpoint customer:
\begin{center}
\texttt{/customer/chatbot/reply/}
\end{center}
Sau đó customer\_service gọi sang chatbot\_service:
\begin{center}
\texttt{http://chatbot-service:8000/api/chat/reply/}
\end{center}

\subsection{Trích đoạn code proxy request}
\begin{lstlisting}[style=pythonstyle,caption={Customer service request chatbot response},label={lst:customer-proxy}]
def request_chatbot_reply(question, current_product=None, user_context=None, user_ref="", limit=5):
    base_url = (os.getenv("CHATBOT_SERVICE_URL") or "http://chatbot-service:8000").rstrip("/")
    endpoint = f"{base_url}/api/chat/reply/"
    payload = {
        "message": question,
        "current_product": current_product,
        "user_context": user_context,
        "user_ref": user_ref,
        "limit": limit,
    }
    response = requests.post(endpoint, json=payload, timeout=40)
    ...
\end{lstlisting}

\subsection{Pipeline command vận hành}
\begin{lstlisting}[style=pythonstyle,caption={Chuỗi command deploy-build-train},label={lst:ops-cmd}]
# build services

docker compose up -d --build chatbot_service customer_service

# backfill du lieu hanh vi thuc te

docker compose exec customer_service python manage.py backfill_chatbot_behavior \
  --max-events 1200 --max-users 300 --source-status paid

# train model_behavior

docker compose exec chatbot_service python manage.py train_behavior_model --epochs 120 --lr 0.02
\end{lstlisting}

\subsection{Vấn đề persistence và cách xử lý}
Trong cấu hình hiện tại, chatbot SQLite nằm trong container filesystem. Khi rebuild image/container, DB có thể reset. Các giải pháp:
\begin{itemize}
  \item Mount volume cho db.sqlite3 và artifacts.
  \item Hoặc tạo quy trình post-rebuild: backfill + train + copy artifacts.
\end{itemize}

\screenshotplaceholder{Screenshot kết quả train + file training\_data\_behavior.json trên VS Code}{fig:training-json}

\section{Đánh giá kết quả theo yêu cầu đề bài}

\subsection{Đối chiếu từng yêu cầu}
\begin{longtable}{p{0.6cm}p{7.0cm}p{2.2cm}p{5.0cm}}
\toprule
STT & Yêu cầu & Mức độ & Bằng chứng \\
\midrule
1 & Khảo sát các áp dụng AI trong e-commerce (5 trang) & Đạt & Mục 2 của báo cáo (tổng quan, bài toán, khoảng trống, tác động KPI) \\
2 & Xây dựng ứng dụng phân tích hành vi để tư vấn & Đạt & BehaviorEvent, ghi event, gateway-request, pipeline chat \\
3 & Xây dựng model\_behavior dựa trên Deep Learning & Đạt & MLP 12-18-10-4, BCE+CE, train command, metrics \\
4 & Xây dựng KB cho tư vấn & Đạt & build\_chat\_kb, knowledge\_base.json, tokenized docs \\
5 & Áp dụng RAG cho chat tư vấn & Đạt & retrieve\_rag\_context, citations, prompt + fallback \\
6 & Deploy và tích hợp trong hệ e-commerce & Đạt & docker compose, inter-service URL, endpoint customer/chatbot \\
\bottomrule
\end{longtable}

\subsection{Chỉ số thực tế sau khi thực thi}
\begin{itemize}
  \item Số event hành vi đã backfill: 1200
  \item Số mẫu train sau augmentation: 3600
  \item Loss lần train gần nhất: 0.009914
  \item File train dễ đọc: training\_data\_behavior.json
\end{itemize}

\section{Hạn chế và hướng phát triển}

\subsection{Hạn chế hiện tại}
\begin{itemize}
  \item Label intent/category hiện tại chủ yếu dựa trên weak labeling và heuristic signal.
  \item Chưa có online A/B test cho recommendation quality.
  \item Chưa tách volume persistence cho chatbot DB trong compose mặc định.
\end{itemize}

\subsection{Hướng nâng cấp}
\begin{itemize}
  \item Bổ sung event-level tracking recommendation\_click, add\_to\_cart, order\_paid theo session timeline.
  \item Chuyển sang train pipeline có train/val/test split theo thời gian.
  \item Bổ sung dashboard theo dõi drift của feature và model.
  \item Nâng cấp RAG retrieval bằng semantic embedding/hybrid search.
\end{itemize}

\section{Kết luận}
Báo cáo đã trình bày đầy đủ toàn bộ yêu cầu bài tập theo hướng triển khai thực tế trên hệ thống e-commerce đã có. Điểm nổi bật là chatbot\_service được tách riêng, kết hợp behavior model + KB + RAG + fallback, đồng thời có quy trình tạo dữ liệu train thực tế và artifact để kiểm tra minh bạch.

Kết quả cho thấy bài toán không chỉ dừng ở mức demo ý tưởng, mà đã được đóng gói thành pipeline vận hành có thể lặp lại: thu thập dữ liệu -> train -> lưu artifact -> tích hợp chat tư vấn vào hệ thống mua bán.

\newpage
\section*{Phụ lục A: Placeholder screenshot cần chèn}
\addcontentsline{toc}{section}{Phụ lục A: Placeholder screenshot cần chèn}

\begin{enumerate}
  \item Tổng quan architecture AI trong e-commerce.
  \item Trạng thái container (docker compose ps).
  \item Kết quả command backfill behavior.
  \item Kết quả train model\_behavior (samples/loss).
  \item Nội dung file knowledge\_base.json.
  \item Giao diện chat có recommendation và citations.
  \item File training\_data\_behavior.json mở trong VS Code.
\end{enumerate}

\section*{Phụ lục B: Hướng dẫn biên dịch LaTeX}
\addcontentsline{toc}{section}{Phụ lục B: Hướng dẫn biên dịch LaTeX}

Có thể biên dịch theo 1 trong 2 cách sau:
\begin{lstlisting}[style=pythonstyle,caption={Build PDF bằng XeLaTeX}]
xelatex REPORT-AI-ECOMMERCE-RAG-DETAILED.tex
xelatex REPORT-AI-ECOMMERCE-RAG-DETAILED.tex
\end{lstlisting}

Hoặc (khi môi trường không có XeLaTeX) biên dịch bằng pdfLaTeX:
\begin{lstlisting}[style=pythonstyle,caption={Build PDF bằng pdfLaTeX}]
pdflatex REPORT-AI-ECOMMERCE-RAG-DETAILED.tex
pdflatex REPORT-AI-ECOMMERCE-RAG-DETAILED.tex
\end{lstlisting}

Nếu bạn sử dụng VS Code LaTeX Workshop, hãy chọn recipe phù hợp:
\begin{itemize}
  \item \texttt{latexmk (xelatex)} khi cần font hệ thống.
  \item \texttt{latexmk (pdflatex)} khi môi trường không có XeLaTeX.
\end{itemize}

Khuyến nghị ưu tiên XeLaTeX nếu máy đã cài sẵn vì cho chất lượng font tốt hơn.

\end{document}
